\par\SourceLeaf{148}\noindent
From the earlier experience, however [p. 139, bottom], taking into account that in its first decades the society was much smaller than in 1794, and that this experience belonged to a time when the corresponding steady state must be regarded as not yet fully reached, one could have concluded that a number of widows far exceeding the average of 13 must henceforth be considered quite possible. (It must not be overlooked that nothing I have said here yet includes the further increase of \(\tfrac16\) required by \textsc{Kritter}'s opinion on account of the orphans' pensions.)

Experience has indeed sufficiently refuted \textsc{Kritter}'s assertion criticized on p. [140], no. 2. It is nevertheless useful to trace the error back to its actual, somewhat hidden source. In every application of mathematical calculation, both to natural phenomena and to social conditions, empirical data seldom occur in the pure form in which they are actually needed. They are almost always affected to a greater or lesser degree by disturbances or fluctuations whose changes follow no rule. As everyone knows, one then tries to remove the resulting disadvantage, or at least reduce it as far as possible, by taking the mean of many individual results. One counts on the random fluctuations largely compensating for one another when a large number of cases is used in this way, and assigns greater reliability to the mean value the more individual results are included. This is entirely correct in general, and the consistent further development and careful use of this principle have often yielded the most rewarding fruits, even brilliant results, especially in the natural sciences. But the soundness of the basic principle rests on an essential condition, often enough overlooked even by specialists: the irregular disturbances or fluctuations affecting the individual observations or experiences must be wholly independent of one another. Judging whether such independence exists can sometimes be very difficult, and impossible without a deep understanding of the circumstances. If doubts remain on this point, the weight to be assigned to the final results will also be uncertain.

Suppose, for example, that we were considering a meteorological quantity, such as the amount of rain falling annually at a particular place. This is known to differ greatly from year to year. But the normal value determined by the general local conditions of the place will be recognized with much greater certainty from an average of ten years than if one relied on a single year. The reason is that there is no particular connection between the deviations from the normal value occurring in the individual years. Rather, as experience also confirms, a large negative deviation can occur just as readily in a year immediately following one with a large positive deviation as in any other year.

\Needspace{32\baselineskip}
But this essential condition is absent in the counts of widows in successive years, precisely because a transition from one number to a substantially different one can occur only gradually. Suppose, for example, that in the larger society used above for illustration, where the \textit{average} annual increase is assumed to be 69 and the annual decrease is equally large once the steady state has been reached, the number has fallen to 1240, or that the negative deviation \(-60\) currently occurs. There is then the greatest improbability, bordering on impossibility, that a positive deviation from the normal value will occur in the following year. In a small society such as ours, the widows counted in the following year are very often still exactly the same individuals as in the preceding year, and even after 10 years only the smaller part will generally have been replaced. An average of 10 successive years is therefore not yet a mean of 10 wholly inde%
\SourceLeaf{149}%
pendent observations, and cannot give much more certain information about the actual normal number than the experience of a single year. To estimate the increased weight to be assigned to an average value, it is essential to consider how many wholly independent groups the observations come from. Now, since the average duration of widowhood is about 20 years, even a forty-year average would still provide no very secure guarantee that the fluctuations had been eliminated when determining the normal number empirically. That for forty successive years the number of widows should remain constantly below the normal value corresponding to the society's general circumstances is, at bottom, no more to be regarded as an altogether extraordinary occurrence than a faro player winning twice in succession. The lesson to be drawn is that, on the other side, in this respect 40 uninterrupted lean years are just as readily possible as 40 uninterrupted fat years.

